\scanpage{12}
\margin{$\overline{\text{\uline{XII}}}$}
\uline{Cor} 5.8. Supposons $S$ \uline{normal} \deleted{. Soit} \inserted{\uline{noeth}.} irréductible. Soit $U$ un ouvert de $S$, \unsure{et} tel que $S-U$ de codim $\geq2$. Alors tt \deleted{tore sur} ss-groupe de t. m. de $G\mid U$ provient d'un unique sous-gpe de t. m. de $G$.

\uline{Cor.} 5.9. Supposons $S$ normal noeth irréd. de pt gén $g$, soit $\deleted{\unsure{H|\overline H}}T_g$ un \deleted{sous-groupe de t. m.}\inserted{sous-tore} de $G_g$. Pour que $H_g$ soit la fibre gén. d'un sous-\deleted{gpe de \unsure{t. m.}}\inserted{tore} $\deleted{\unsure{H|\overline H}}T$ de $G$, il f. et s. que pour tt $s\in S$ tel que $\deleted{\dim(\unsure{H|T|\overline H}_g)}$ $\dim\unsure{\mathcal O|\mathscr O}_{S,s}=1$, $\left(\overline{H_g}\right)^{\circ}_{s,\deleted{\unsure{\mathrm{réd}}}\mathrm{réd}}$ soit \deleted{un tore.} de t. m.
% slanted
\margin{\uline{rédiger} \uline{mieux}}

\deleted{\uline{Cor.} 5.10. Sous les conditions précédentes, supposons que $H_g$ soit $=$ un \uline{tore} [plus gén., que \deleted{\unsure{son type n'ait pas de}}\inserted{les \unsure{diviseurs} de la} \unsure{p-torsion|torsion} de son type soient $\unsure{\ne|\pm}$ car. résiduelles] \uline{Alors}}

On utilise un

\uline{Lemme} \ldots

\margin{TSVP}
% archive sheet number: 11
